% Generated draft: requires mathematical and source review.
\documentclass[11pt]{article}
\usepackage[margin=25mm]{geometry}
\usepackage{fontspec}
\setmainfont{Latin Modern Roman}
\usepackage{amsmath,amssymb,mathtools,unicode-math}
\setmathfont{Latin Modern Math}
\usepackage{enumitem,needspace,xurl,hyperref}
\setlength{\parindent}{0pt}
\setlength{\parskip}{5pt}
\setlength{\emergencystretch}{3em}
\newcommand{\N}{\mathbb{N}}
\newcommand{\Z}{\mathbb{Z}}
\newcommand{\Q}{\mathbb{Q}}
\newcommand{\R}{\mathbb{R}}
\newcommand{\C}{\mathbb{C}}
\newcommand{\F}{\mathbb{F}}
\newcommand{\A}{\mathbb{A}}
\newcommand{\PP}{\mathbb{P}}
\newcommand{\E}{\mathbb{E}}
\newcommand{\eps}{\varepsilon}
\newcommand{\dd}{\,\mathrm{d}}
\begin{document}
\section*{UnsolvedMath ID: 10000009}
\textbf{Source problem code:} catalogue entry
\subsection{Definitions and mathematical statement}
A graph is vertex transitive if its automorphism group acts transitively on its vertices. Equip it with the graph path metric. A common $k$-net is a vertex subset that is $k$-separated and coarsely covers the graph at radius comparable to $k$, but the source does not define its exact separation/covering convention or the edges and edge lengths of the associated net graph.

For specified graph and net metrics, a rough-isometry comparison with multiplicative constant $\lambda_k\ge1$ would use a map $f_k$ and additive control $A_k\ge0$ satisfying
\[
 \lambda_k^{-1}d_G(x,y)-A_k\le
 d_{N_k}(f_k(x),f_k(y))\le\lambda_k d_G(x,y)+A_k,
\]
with every point of $N_k$ within the allowed coarse-surjectivity distance of the image. The source asks whether there exists a vertex-transitive graph for which the relevant multiplicative constants can satisfy
\[
 \lambda_k\longrightarrow1\qquad(k\to\infty).
\]
It does not prescribe whether additive and coarse-surjectivity constants must be uniform in $k$, whether distances are rescaled at net scale, or whether the constants are optimized over maps.

The net quantifier is also missing: a chosen net at each scale, every admissible net, and nets selected by an invariant rule are different requirements. Connectedness, local finiteness, infinitude, and nontriviality assumptions on the graph are not explicitly given. These conventions are essential to the intended sub-scale-invariance problem and its possible degenerate cases, so no particular net construction or simple example is silently chosen.
\subsection{Short English statement}
Can a vertex-transitive graph become arbitrarily close to multiplicative distortion one when compared with its coarse net graphs? The net construction, rescaling, and additive-control requirements must be fixed.

The supplied summary calls the problem open and mentions scale-invariance/asymptotic-dimension context without exact references. It calls the notes 2024, while provenance includes a 2020 archived copy; the document dating is not inferred from the filename. The recovered-source and bounded-modern-audit assessments are retained without an independent current-status claim.
\subsection{Sources}
\begin{itemize}
\item Local UnsolvedMath registry, permanent ID 10000009, source code catalogue entry, ``Vertex-transitive sub-scale-invariant graphs,'' proposed by Itai Benjamini. Dataset curation and metadata: Ulam AI and UnsolvedMath Contributors, CC BY 4.0.
\item Itai Benjamini, \emph{Coarse Geometry and Randomness}, Saint-Flour notes, Open problem 1.67, PDF page 14, \url{https://www.wisdom.weizmann.ac.il/~itai/stflouraug24.pdf#page=14}, preserved at \url{https://arquivo.pt/noFrame/replay/20201231041548id_/http://www.wisdom.weizmann.ac.il/~itai/stflouraug24.pdf}; accessed 2026-07-29, source extraction reported as pypdf text checked against rendered PDF.
\item No exact related scale-invariance theorem is supplied. Provenance says cross-document deduplication and source-side status review retained the item, with a bounded modern audit finding no complete primary resolution.
\item Source-specific current status and release reuse terms NEEDS\_REVIEW, separately from dataset attribution. No source visit or audit is claimed in this draft.
\end{itemize}
\end{document}
